\documentclass[11pt,a4paper]{article}
\usepackage[slovak]{babel}
\usepackage[utf8]{inputenc}
\usepackage[T1]{fontenc}
\usepackage{geometry}
\geometry{margin=2.5cm}
\usepackage{graphicx}
\usepackage{float}
\usepackage{hyperref}

\title{TimeManager (TIMA): jednoduchý plánovač času pre študentov}
\author{Michal Pípa, Tomáš Topalov, Tomáš Šuba, Michal Urban}
\date{Akademický rok 2025/26}

\begin{document}
\maketitle

\begin{abstract}
Článok opisuje návrh aplikácie TimeManager (TIMA), jednoduchej pomôcky pre efektívne plánovanie času študentov.
TIMA využíva princípy time managementu, ako sú Eisenhowerova matica a Pomodoro technika, a umožňuje používateľom
vytvárať, triediť a plánovať úlohy s pripomienkami. Projekt je zameraný na jednoduchosť a praktickú využiteľnosť.
\end{abstract}

\section{Úvod}
Študenti často zápasia s organizáciou svojho času. TIMA pomáha jednoducho plánovať úlohy, nastavovať priority a kontrolovať pokrok.

\section{Ciele}
\begin{itemize}
\item umožniť efektívne plánovanie školských úloh,
\item znížiť stres z neorganizovaného rozvrhu,
\item podporiť disciplínu pomocou pripomienok.
\end{itemize}

\section{Návrh riešenia}
Aplikácia obsahuje tri hlavné časti:
\begin{itemize}
\item \textbf{Inbox} – rýchle pridávanie úloh,
\item \textbf{Prioritizácia} – zoradenie podľa dôležitosti (Eisenhowerova matica),
\item \textbf{Time-blocking} – plánovanie úloh do časového rozvrhu.
\end{itemize}

\section{Tabuľka úloh}
\begin{table}[H]
\centering
\begin{tabular}{|l|c|c|c|}
\hline
Úloha & Odhad [h] & Priorita (1--4) & Skóre = Odhad × Priorita \\
\hline
Matematika & 2 & 3 & 6 \\
Projekt MIP & 1 & 4 & 4 \\
Angličtina & 1.5 & 2 & 3 \\
\hline
\end{tabular}
\caption{Ukážka tabuľky úloh v systéme TIMA.}
\end{table}

\section{Záver}
TIMA predstavuje jednoduchý, prehľadný a použiteľný nástroj, ktorý môže pomôcť študentom efektívne riadiť čas a zlepšiť produktivitu.

\bibliographystyle{plain}
\bibliography{refs_tima}
\end{document}